% ===================== ANALISIS DE RESULTADOS =====================
\section{Analisis de Resultados}
\label{sec:analisis}

% ─────────────────────────────────────────────────────────────────────────────
\subsection{Cumplimiento de Codigo (Code Compliance)}

La evaluaci\'{o}n de cumplimiento de c\'{o}digo del modelo arroj\'{o}
\textbf{CODE COMPLIANCE EVALUATION PASSED}. El esfuerzo de c\'{o}digo
m\'{a}ximo del sistema es \num{39668.2}~\si{kPa}, en el nodo 78, para el caso
6 (SUS) W+P2, frente a un admisible de \num{115142.4}~\si{kPa}; la
relaci\'{o}n de utilizaci\'{o}n m\'{a}xima es por tanto 34.5~\%. La
Tabla~\ref{tab:top5-ratio} relaciona los cinco nodos de mayor ratio del caso
cr\'{i}tico; todos cumplen con holgura el criterio Code $\leq$ Allowable.

\begin{table}[htbp]
  \centering
  \small
  \caption{Nodos de mayor ratio de esfuerzo — línea SIM-002. Criterio: Code $\leq$ Allowable.}
  \label{tab:top5-ratio}
  \begin{tabular}{crrrrrc}
    \toprule
    N & Nodo & Bending (\si{kPa}) & Code (\si{kPa}) & Allowable (\si{kPa}) & Ratio (\%) & Cumple \\
    \midrule
    1 & 78 & \num{26891.5} & \num{39668.2} & \num{115142.4} & 34.5 & Sí \\
    2 & 90 & \num{21279.9} & \num{34805.8} & \num{115142.4} & 30.2 & Sí \\
    3 & 90 & \num{21279.9} & \num{33350.8} & \num{115142.4} & 29.0 & Sí \\
    4 & 70 & \num{11959.7} & \num{24736.5} & \num{115142.4} & 21.5 & Sí \\
    5 & 70 & \num{11959.7} & \num{24736.5} & \num{115142.4} & 21.5 & Sí \\
    \bottomrule
  \end{tabular}
\end{table}


El caso que gobierna el dise\~{n}o es el sostenido W+P2, es decir, la
combinaci\'{o}n de peso propio con la presi\'{o}n de operaci\'{o}n $P_2$. Este
resultado es caracter\'{i}stico de l\'{i}neas cortas y r\'{i}gidas entre
restricciones, donde el esfuerzo longitudinal por presi\'{o}n y peso domina
sobre el rango t\'{e}rmico: los casos de expansi\'{o}n presentan ratios
moderados (m\'{a}ximo 28.0~\% en el nodo 90, caso 8), frente a un admisible de
expansi\'{o}n de \num{254505.3}~\si{kPa}. El esfuerzo de flexi\'{o}n se
concentra en la zona de los nodos 70 a 90, como evidencia el perfil de la
Figura~\ref{fig:esfuerzos}, con el m\'{a}ximo en el nodo 78
(\num{26891.5}~\si{kPa} de flexi\'{o}n en el caso 6). En condici\'{o}n de
operaci\'{o}n (caso 3, OPE, sin verificaci\'{o}n de c\'{o}digo) el modelo
reporta un esfuerzo de flexi\'{o}n m\'{a}ximo de \num{84639.1}~\si{kPa} en el
nodo 90 y una intensidad de esfuerzo m\'{a}xima de \num{100019.5}~\si{kPa} en
el mismo nodo.

% ─────────────────────────────────────────────────────────────────────────────
\subsection{Desplazamientos y Cargas en Soportes}

Los desplazamientos de operaci\'{o}n est\'{a}n controlados por las holguras de
los soportes de los extremos: el cierre de la holgura limita el movimiento
horizontal a 5.0~\si{mm} y los desplazamientos verticales no superan
2.0~\si{mm} (Tabla~\ref{tab:desplaz-max}), lo que indica que la l\'{i}nea
trabaja con bajos movimientos relativos y sin riesgo de interferencias por
expansi\'{o}n.

Las cargas en soportes (Tablas~\ref{tab:restricciones-ope} a
\ref{tab:restricciones-exp}) son moderadas: la carga vertical m\'{a}xima
corresponde al soporte del nodo 70 con \num{-5529}~\si{N} en el caso 1 (OPE) y
\num{-9013}~\si{N} en el caso 3 (OPE con $T_2$), y las restricciones de los
extremos transmiten fuerzas reducidas una vez cerrada la holgura. No obstante,
la comparaci\'{o}n de las cargas en la boquilla de la bomba con los admisibles
del fabricante del equipo (p.\,ej. API 610) queda fuera del alcance del
presente informe (secci\'{o}n~\ref{sec:alcance}).

% ─────────────────────────────────────────────────────────────────────────────
\subsection{Sintesis de Hallazgos}

La l\'{i}nea SIM-002 presenta un margen amplio frente a los criterios de
ASME B31.3-2016: el ratio m\'{a}ximo de utilizaci\'{o}n (34.5~\%) deja cerca
de dos tercios del admisible disponible, los desplazamientos son peque\~{n}os
(m\'{a}ximo 5.0~\si{mm}, controlados por el cierre de holguras) y las cargas en
soportes son moderadas. No se identifican cuellos de botella, zonas de
sobre-esfuerzo ni configuraciones de soporte cr\'{i}ticas. El \'{u}nico aspecto
pendiente de verificaci\'{o}n externa es la aceptaci\'{o}n de las cargas en la
boquilla de la bomba por parte de su fabricante.
